\section{Fundamentação teórica e trabalhos relacionados}
\label{sec:fundamentacao}

\subsection{Jogos sérios e educacionais}

O termo \emph{jogos sérios} popularizou-se a partir da definição de Abt:
jogos com um propósito educacional explícito, que não são feitos apenas para
divertir --- embora possam ser divertidos \citep{abt1970serious}. A pesquisa em
aprendizagem com jogos mostrou que bons jogos já incorporam, de forma orgânica,
princípios que a escola persegue: identidade, experimentação, risco controlado,
regras claras e retroalimentação constante \citep{gee2003what}. A visão de
simulação virtual como base para jogos, e dos jogos como ferramenta de
treinamento, foi consolidada na transição de simulação para realidade
virtual \citep{zyda2005from}. Revisões mais recentes organizam o campo por
aplicação (saúde, defesa, educação, planejamento) e apontam que a avaliação de
impacto segue sendo um desafio metodológico
\citep{laamarti2014overview}. [REFERÊNCIA NECESSÁRIA: diretrizes
contemporâneas de avaliação de jogos sérios em contextos escolares.]

\subsection{Grandes modelos de linguagem em jogos}

O levantamento de Gallotta \emph{et al.} mapeia os papéis que LLMs podem
assumir em jogos --- geração de conteúdo, personagens, apoio ao projeto,
testes --- e discute limites práticos, como custo, latência e confiabilidade
\citep{gallotta2024largelanguage}. Duas limitações recorrentes interessam
diretamente a este trabalho: (a) modelos grandes exigem infraestrutura
relevante ou serviços externos; e (b) quanto maior a autonomia concedida ao
modelo, maior a superfície de risco sobre o estado do jogo.

\subsection{Agentes generativos e autonomia}

\emph{Generative Agents} propôs agentes que armazenam experiências em linguagem
natural, sintetizam reflexões e planejam comportamento, produzindo
comportamento social emergente e verossímil em um ambiente interativo
\citep{park2023generative}. \emph{Voyager} levou a autonomia adiante em
Minecraft: um currículo automático, uma biblioteca de habilidades em código e
um mecanismo iterativo de prompting com retroalimentação do ambiente; o agente
aprende continuamente sem ajuste de parâmetros, consultando um modelo externo
\citep{wang2023voyager}. Esses trabalhos demonstram o potencial da autonomia,
mas também a distância entre esse desenho e um jogo educacional em que cada
ação precisa ser previsível, auditável e barata de executar. Neste artigo, a
escolha é deliberadamente oposta: a IA conversa, comenta e sugere; quem age é
sempre o motor do jogo.

\subsection{Inferência local}

Executar o modelo na própria máquina elimina dependência de serviços externos,
dispensa chaves de API e mantém o conteúdo no computador do usuário; o Ollama
é uma ferramenta consolidada para esse cenário \citep{ollama2026}. O custo é a
latência: modelos menores rodam em hardware modesto, mas geram respostas em
escala de segundos, o que restringe seu uso a interações conversacionais (por
turnos) em vez de reações em tempo real. [REFERÊNCIA NECESSÁRIA: estudo
comparativo de latência de LLMs locais em hardware sem GPU dedicada.]

\subsection{Processo de desenvolvimento}

Práticas como desenvolvimento orientado a testes e a especificações propõem
que critérios verificáveis sejam escritos antes da implementação.
[REFERÊNCIA NECESSÁRIA: fontes canônicas sobre desenvolvimento orientado a
especificação.] No contexto deste projeto, essas práticas foram adaptadas a um
fluxo de trabalho assistido por agentes de IA: a especificação, os requisitos e
os critérios de aceitação são os artefatos que restringem cada iteração de
implementação, e a suíte automatizada é o juiz comum entre autor humano e
agente. O uso de servidores MCP (\emph{Model Context Protocol}) para dar ao
agente acesso controlado a ferramentas --- no caso, a cena do Blender --- é
parte desse fluxo \citep{lab404repo}.

\subsection{Síntese e lacuna}

A literatura concentra-se, de um lado, em agentes autônomos e capazes de agir
\citep{park2023generative,wang2023voyager} e, de outro, em visões amplas sobre
o papel das LLMs em jogos \citep{gallotta2024largelanguage}. O que este
trabalho investiga é a combinação menos explorada: \emph{uma IA diegética, com
capacidade formalmente limitada, executada localmente e integrada a um jogo
educacional completo, com processo de desenvolvimento documentado e
evidências de validação}. É nessa lacuna que se posicionam as contribuições
descritas na \autoref{sec:metodologia}.
